% TeleCon Classic Simple — guía del código del sketch Arduino.
% Build: pdflatex -output-directory=../build telecon_classic_simple_code.tex
% PDF: ../build/telecon_classic_simple_code.pdf
\documentclass[11pt,a4paper]{article}

\usepackage[margin=2.2cm]{geometry}
\usepackage[T1]{fontenc}
\usepackage[utf8]{inputenc}
\usepackage{lmodern}
\usepackage{booktabs}
\usepackage{hyperref}
\usepackage{xcolor}
\usepackage{colortbl}
\usepackage{array}
\usepackage{enumitem}
\usepackage{needspace}
\usepackage{etoolbox}
\usepackage{listings}
\usepackage{upquote}
\usepackage{tikz}
\usepackage[breakable]{tcolorbox}
\usetikzlibrary{arrows.meta,positioning,shapes.geometric,fit,backgrounds,matrix}

\preto\section{\Needspace{8\baselineskip}}
\preto\subsection{\Needspace{5\baselineskip}}

% Wide enough for "Anexo A" in the contents, on the same line as the file name.
\makeatletter
\patchcmd{\l@section}{\setlength\@tempdima{1.5em}}{\setlength\@tempdima{5.8em}}{}{}
\makeatother

\hypersetup{
  colorlinks=true,
  linkcolor=black,
  urlcolor=blue,
  pdfhighlight=/N,
  pdftitle={Guía de código TeleCon Classic Simple},
  pdfauthor={TeleCon}
}

\definecolor{TblHead}{HTML}{0F6F73}
\definecolor{TblChan}{HTML}{D5F5E3}
\definecolor{TblEdit}{HTML}{F5CBA7}
\definecolor{TblTelem}{HTML}{D4E6F1}
\definecolor{TblLeave}{HTML}{F4F6F7}

% Arduino IDE 1.8 editor (theme.txt + keywords.txt).
\definecolor{ArduinoType}{HTML}{00979C}
\definecolor{ArduinoFunc}{HTML}{D35400}
\definecolor{ArduinoControl}{HTML}{5E6D03}
\definecolor{ArduinoString}{HTML}{005C5F}
\definecolor{ArduinoComment}{HTML}{434F54}
\definecolor{ArduinoName}{HTML}{6C3483}
\definecolor{ArduinoLineNo}{HTML}{7F8C8D}
\definecolor{ArduinoFrame}{HTML}{C5C9C9}

\newcommand{\thc}[1]{\textcolor{white}{\bfseries #1}}
\newcommand{\chip}[2]{\colorbox{#1}{\strut\,\scriptsize #2\,}}
\newcommand{\ui}[1]{\textsf{\textbf{#1}}}
\newcommand{\code}[1]{\texttt{#1}}
\newcommand{\nm}[1]{\textcolor{ArduinoName}{\texttt{#1}}}
\newcommand{\fnm}[1]{\textcolor{ArduinoFunc}{\texttt{#1}}}
\newcommand{\typ}[1]{\textcolor{ArduinoType}{\texttt{#1}}}
\newcommand{\kw}[1]{\textcolor{ArduinoControl}{\texttt{#1}}}
\newcommand{\str}[1]{\textcolor{ArduinoString}{\texttt{#1}}}
\newcommand{\cmt}[1]{\textcolor{ArduinoComment}{\texttt{#1}}}
\newcommand{\ser}[1]{\textcolor{ArduinoFunc}{\textbf{\texttt{Serial}}\texttt{.#1}}}

\newcommand{\sketchfile}{../../../ESP32_noCam/Classic_Simple/TeleCon_Classic_Simple/TeleCon_Classic_Simple.ino}
\newcommand{\channelfile}{../../../ESP32_noCam/Classic_Simple/TeleCon_Classic_Simple/TelemetryChannels.h}

\lstdefinelanguage{Arduino}{
  language=C++,
  alsoletter={\_},
  keywordstyle=\color{ArduinoType},
  keywordstyle=[2]\color{ArduinoFunc},
  keywordstyle=[3]\color{ArduinoControl},
  keywordstyle=[4]\color{ArduinoFunc}\bfseries,
  keywordstyle=[5]\color{ArduinoName},
  commentstyle=\color{ArduinoComment},
  stringstyle=\color{ArduinoString},
  directivestyle=\color{ArduinoControl},
  deletekeywords={%
    auto,break,case,char,const,continue,default,do,double,else,enum,%
    extern,float,for,goto,if,int,long,register,return,short,signed,%
    sizeof,static,struct,switch,typedef,union,unsigned,void,volatile,%
    while,and,and_eq,asm,bad_cast,bad_typeid,bitand,bitor,bool,catch,%
    class,compl,const_cast,delete,dynamic_cast,explicit,export,false,%
    friend,inline,mutable,namespace,new,not,not_eq,operator,or,or_eq,%
    private,protected,public,reinterpret_cast,static_cast,template,%
    this,throw,true,try,typeid,type_info,typename,using,virtual,%
    wchar_t,xor,xor_eq,alignas,alignof,char16_t,char32_t,constexpr,%
    decltype,noexcept,nullptr,static_assert,thread_local%
  },
  morekeywords={%
    void,int,char,float,double,long,short,signed,unsigned,bool,boolean,%
    byte,word,const,static,struct,enum,class,sizeof,typedef,volatile,%
    inline,constexpr,String,uint8_t,uint16_t,uint32_t,uint64_t,%
    int8_t,int16_t,int32_t,int64_t,true,false,nullptr,NULL,HIGH,LOW,%
    INPUT,OUTPUT,INPUT_PULLUP%
  },
  morekeywords=[2]{%
    print,println,printf,begin,available,read,readStringUntil,trim,%
    length,indexOf,substring,startsWith,equalsIgnoreCase,toInt,c_str,%
    hasClient,disableSSP,setPin,delay,millis,map,constrain,abs,%
    analogRead,digitalRead,ledcAttach,ledcWrite,ledcRead,ledcSetup,%
    ledcAttachPin,snprintf,strtol,%
    setup,loop,%
    telemetryVehicle,steerDutyFromUs,steerDutyFromRx,%
    steerServoBegin,steerServoWriteRx,onControl,onButton,%
    handleSerialInjectLine,serialInjectPoll,stickTo12,%
    debugPrintControl,debugPrintButton,sendLine,applyFailSafe,%
    resetHandshake,getValue,handleConnect,handleLine,%
    telemetryReadPwm,telemetryReadOne,telemetryReadAll%
  },
  morekeywords=[3]{%
    if,else,for,while,do,switch,case,default,break,continue,return,%
    goto%
  },
  morekeywords=[4]{Serial},
  morekeywords=[5]{%
    BT_NAME,BT_PIN,APP_PREFIX,PROTO_WIRE,SERIAL_BAUD,%
    ENABLE_DEBUG,DBG_PRINT,DBG_PRINTLN,DBG_PRINTF,%
    TELEMETRY_DATA_MS,ENABLE_PLOT,PLOT_MS,%
    ENABLE_SERIAL_INJECT,SERIAL_MAX_LINE,%
    DEMO_MOTOR_TEMP,DEMO_BATT_PCT,%
    ENABLE_STEER_SERVO,PIN_STEER_SERVO,CH_STEER_SERVO,%
    ESP_ARDUINO_VERSION_MAJOR,%
    SerialBT,handshakeOk,btWasConnected,%
    g_lx,g_ly,g_rx,g_ry,g_lk,g_rk,g_sw,%
    lastDataMs,lastPlotMs,serialLineBuf,%
    lastPrintLX,lastPrintLY,lastPrintRX,lastPrintRY,%
    lastPrintKnobL,lastPrintKnobR,lastPrintSw,%
    kTelemetryChannels,kSteerCenterUs,kSteerTravelUs,%
    TelKind,TelVehicle,TelChannel,TelVehicleFn,%
    TEL_VEHICLE,TEL_ADC,TEL_GPIO,TEL_PWM,%
    TEL_TEMP,TEL_THROTTLE,TEL_BATT,TEL_STEER,TEL_LIGHTS,%
    TEL_KNOB_L,TEL_KNOB_R,TEL_RY,%
    kind,vehicle,pin,ledcChannel,dutyMin,dutyMax%
  },
  literate=
    {µ}{{\textmu}}1
    {×}{{$\times$}}1
    {–}{{\textendash}}1
    {—}{{\textemdash}}1
    {…}{{\ldots}}1
    {→}{{$\rightarrow$}}1
}

\lstset{
  language=Arduino,
  basicstyle=\ttfamily\normalsize,
  numbers=left,
  numberstyle=\footnotesize\color{ArduinoLineNo},
  numbersep=6pt,
  stepnumber=1,
  frame=single,
  framerule=0.4pt,
  rulecolor=\color{ArduinoFrame},
  framesep=3pt,
  framexleftmargin=0pt,
  backgroundcolor=\color{white},
  breaklines=true,
  breakindent=0pt,
  columns=fullflexible,
  keepspaces=true,
  showstringspaces=false,
  showtabs=false,
  tabsize=2,
  upquote=true,
  xleftmargin=2.15em,
  aboveskip=0.65em,
  belowskip=0.45em
}

\newcommand{\inosnippet}[2]{%
  \par\addvspace{1.35em}%
  \lstinputlisting[firstline=#1,lastline=#2,firstnumber=#1,aboveskip=0.4em]{\sketchfile}%
}
\newcommand{\hsnippet}[2]{%
  \par\addvspace{1.35em}%
  \lstinputlisting[firstline=#1,lastline=#2,firstnumber=#1,aboveskip=0.4em]{\channelfile}%
}
\newcommand{\fnat}[4]{%
  \par\addvspace{1.15em}%
  \noindent{\ttfamily\Large\raggedright #1\par}%
  \nopagebreak
  \lstinputlisting[firstline=#3,lastline=#4,firstnumber=#3,%
    aboveskip=0.55em,belowskip=0.35em,%
    backgroundcolor=\color{FnShade}]{#2}%
  \noindent\ignorespaces
}
\newcommand{\fn}[3]{\fnat{#1}{\sketchfile}{#2}{#3}}
\newcommand{\fnh}[3]{\fnat{#1}{\channelfile}{#2}{#3}}
\newcommand{\fntwo}[5]{%
  \par\addvspace{1.15em}%
  \noindent{\ttfamily\Large\raggedright #1\par}%
  \nopagebreak
  \lstinputlisting[firstline=#2,lastline=#3,firstnumber=#2,%
    aboveskip=0.55em,belowskip=0.2em,%
    backgroundcolor=\color{FnShade}]{\sketchfile}%
  \lstinputlisting[firstline=#4,lastline=#5,firstnumber=#4,%
    aboveskip=0.2em,belowskip=0.35em,%
    backgroundcolor=\color{FnShade}]{\sketchfile}%
  \noindent\ignorespaces
}
% Snippet on the left, title and explanation on the right.
% Measured boxes so the code frame sits on the vertical center of the text.
\newsavebox{\fnstepcode}
\newsavebox{\fnsteptext}
\newcommand{\fnstepat}[5]{%
  \par\addvspace{1.7em}%
  \sbox{\fnstepcode}{%
    \begin{minipage}[b]{0.50\linewidth}
      \lstinputlisting[firstline=#2,lastline=#3,firstnumber=#2,%
        basicstyle=\ttfamily\footnotesize,%
        numberstyle=\scriptsize\color{ArduinoLineNo},%
        numbersep=5pt,framesep=3pt,framexleftmargin=0pt,%
        xleftmargin=2.2em,aboveskip=0pt,belowskip=0pt,%
        breaklines=true,breakindent=0pt,%
        backgroundcolor=\color{FnShade}]{#1}%
    \end{minipage}}%
  \sbox{\fnsteptext}{%
    \begin{minipage}[b]{0.44\linewidth}
      {\raggedright\Large\bfseries #4\par}%
      \vspace{0.45em}%
      \raggedright #5\par
    \end{minipage}}%
  \noindent
  % Listings' frame sits a little above the measured box; nudge the code down.
  \raisebox{\dimexpr(\ht\fnsteptext-\ht\fnstepcode+\dp\fnstepcode-\dp\fnsteptext)/2-8pt\relax}{\usebox{\fnstepcode}}%
  \hspace{0.04\linewidth}%
  \usebox{\fnsteptext}%
  \par
}
\newcommand{\fnstep}[4]{\fnstepat{\sketchfile}{#1}{#2}{#3}{#4}}
\newcommand{\fnhstep}[4]{\fnstepat{\channelfile}{#1}{#2}{#3}{#4}}

\definecolor{FnShade}{gray}{1}
\newcounter{fnband}
\newenvironment{fnblock}{%
  \par
  \stepcounter{fnband}%
  \ifodd\value{fnband}%
    \colorlet{FnShade}{black!4}%
  \else
    \colorlet{FnShade}{white}%
  \fi
  \begin{tcolorbox}[%
    breakable,colback=FnShade,colframe=FnShade,boxrule=0pt,%
    arc=0pt,left=18pt,right=18pt,top=6pt,bottom=8pt,boxsep=0pt,%
    before skip=0.6em,after skip=1.15em,pad at break*=2pt]%
}{%
  \end{tcolorbox}%
}

\tikzset{
  flow/.style={draw=TblHead, fill=TblChan, rounded corners=2pt,
               align=center, font=\sffamily\small, inner sep=5pt,
               minimum height=8mm},
  edit/.style={draw=TblHead, fill=TblEdit, rounded corners=2pt,
               align=center, font=\sffamily\small\bfseries, inner sep=5pt},
  leave/.style={draw=TblHead!70, fill=TblLeave, rounded corners=2pt,
                align=center, font=\sffamily\small, inner sep=5pt},
  arr/.style={-{Stealth[length=2.2mm]}, draw=TblHead, thick}
}

\title{Classic + Simple\\[0.25em]
\large Guía de código de \texttt{TeleCon\_Classic\_Simple}}
\author{ESP32 DevKit \textperiodcentered\ Bluetooth Classic \textperiodcentered\ líneas de texto}
\date{Versión 1.0 \quad 2026}

\begin{document}
\maketitle

\begin{abstract}
Esta nota es el sketch de Arduino en la carpeta
\texttt{Classic\_Simple/TeleCon\_Classic\_Simple/}.
El ajuste del teléfono es Classic + Simple. El nombre de la placa es
\str{ESP32-TC-RC-BT-Simple}. El saludo es
\str{RC:CONNECT,proto,simple}.
Primero van las declaraciones de cada archivo: includes, macros,
variables y tipos. Después cada función, con su código en los colores
del IDE.
Los dos archivos completos quedan como anexos.
\end{abstract}

{\sffamily\scriptsize
Color del IDE:
\textcolor{ArduinoType}{\texttt{void}},
\textcolor{ArduinoControl}{\texttt{if}},
\textcolor{ArduinoControl}{\texttt{\#define}},
\textcolor{ArduinoFunc}{\textbf{\texttt{Serial}}}\textcolor{ArduinoFunc}{\texttt{.print}},
\textcolor{ArduinoString}{\texttt{"texto"}},
\textcolor{ArduinoName}{\texttt{BT\_NAME}},
\textcolor{ArduinoComment}{\texttt{// comentario}}.}

\tableofcontents
\clearpage

\section{Declaraciones de los archivos}
\label{sec:decl}

Nada de esto es una función. Es la cabecera de cada archivo: lo que
el compilador ve antes del primer cuerpo. El código va con el color
del IDE. Los números de línea son los del archivo.

\subsection{En \texttt{TelemetryChannels.h}}

El \texttt{.ino} incluye esta cabecera, así que estos nombres existen
antes de \nm{kTelemetryChannels[]}.

\hsnippet{1}{15}

\kw{\#pragma once} hace que el archivo entre una sola vez en la
compilación, aunque alguien lo incluya dos veces. \texttt{Arduino.h}
trae \typ{uint8\_t}, \fnm{analogRead}, \fnm{digitalRead} y
\fnm{ledcRead}. El comentario de en medio no genera código: recuerda
que CH1 es \str{v0}, CH8 es \str{v7}, y que cada muestra es un
byte 0--255.

\hsnippet{17}{33}

\nm{TelKind} dice de dónde sale el byte de una fila.
\nm{TEL\_VEHICLE} llama a \fnm{telemetryVehicle()}.
\nm{TEL\_ADC} lee un pin analógico. \nm{TEL\_GPIO} lee un pin
digital. \nm{TEL\_PWM} lee un duty LEDC. El valor numérico empieza
en 0 y sigue en orden.

\nm{TelVehicle} nombra la señal integrada, solo cuando la fila es
\nm{TEL\_VEHICLE}. \nm{TEL\_TEMP} es temperatura,
\nm{TEL\_THROTTLE} el acelerador, \nm{TEL\_BATT} la batería,
\nm{TEL\_STEER} la dirección, \nm{TEL\_LIGHTS} las luces,
\nm{TEL\_KNOB\_L} y \nm{TEL\_KNOB\_R} las perillas, y
\nm{TEL\_RY} el eje Y de la palanca derecha. La tabla de fábrica
no usa \nm{TEL\_RY}.

\hsnippet{35}{44}

\nm{TelChannel} es una fila de la tabla. \nm{kind} es el
\nm{TelKind}. \nm{vehicle} es el \nm{TelVehicle}; en una
fila ADC, GPIO o PWM ese campo no se lee. \nm{pin} es el GPIO, o
$-1$ si la fila no tiene pin. \nm{ledcChannel} es el canal LEDC
del núcleo 2.x, o $-1$ si no hay PWM. \nm{dutyMin} y
\nm{dutyMax} son los extremos del duty: el mínimo se manda como 0
y el máximo como 255.

\nm{TelVehicleFn} es el tipo del puntero a
\fnm{telemetryVehicle()}: recibe un \nm{TelVehicle} y devuelve
un byte.

\subsection{En \texttt{TeleCon\_Classic\_Simple.ino}}

\inosnippet{1}{14}

El bloque de comentario no entra al binario. Dice el nombre Bluetooth,
el PIN y la línea de saludo. \texttt{BluetoothSerial.h} es el puerto
Classic. \texttt{TelemetryChannels.h} es la cabecera de arriba.

\inosnippet{16}{45}

\nm{BT\_NAME} es el nombre que Android ve al emparejar.
\nm{BT\_PIN} es \str{1234}. \nm{APP\_PREFIX} tiene que ser
\str{RC} y \nm{PROTO\_WIRE} tiene que ser \str{simple}: son
las dos cadenas que acepta \fnm{handleConnect()}.
\nm{SERIAL\_BAUD} es 115200 para el monitor.

\nm{ENABLE\_DEBUG} en 1 escribe en el monitor.
\nm{DBG\_PRINT} pasa a \ser{print}.
\nm{DBG\_PRINTLN} pasa a \ser{println}.
\nm{DBG\_PRINTF} pasa a \ser{printf}.
En 0 esas tres macros quedan en nada y el enlace sigue igual.

\nm{TELEMETRY\_DATA\_MS} es el periodo del tablero, 500\,ms.
\nm{ENABLE\_PLOT} en 1 envía \str{RC:PLOT}.
\nm{PLOT\_MS} es 50\,ms.
\nm{ENABLE\_SERIAL\_INJECT} en 1 deja pegar una línea en el monitor
después del saludo.
\nm{SERIAL\_MAX\_LINE} corta esa línea a 192 caracteres.

\nm{DEMO\_MOTOR\_TEMP} es 90 y \nm{DEMO\_BATT\_PCT} es 76.
Son números fijos del tablero hasta que haya un sensor.
No salen de un pin.

\inosnippet{47}{70}

\nm{SerialBT} es el puerto Classic. \nm{handshakeOk} es la
compuerta: en falso no hay \fnm{onControl()}, ni \str{RC:DATA},
ni \str{RC:PLOT}. \nm{btWasConnected} guarda el cliente de la
vuelta anterior, para notar una conexión o una caída una sola vez.

\nm{g\_lx}, \nm{g\_ly}, \nm{g\_rx} y \nm{g\_ry} son las
palancas, de $-100$ a 100. Arrancan en 0. \nm{g\_ly} es el
acelerador y \nm{g\_rx} la dirección. \nm{g\_lk} y
\nm{g\_rk} son las perillas, de 0 a 1023, y arrancan en 512.
\nm{g\_sw} es el byte de interruptores. El bit 0 son las luces.

\nm{lastDataMs} y \nm{lastPlotMs} son las marcas de
\fnm{millis()} del tablero y del gráfico. \nm{serialLineBuf}
junta la línea que se pega en el monitor; solo existe si la inyección
está compilada. Las \nm{lastPrint*} guardan la última imagen
impresa de cada eje, perilla e interruptor. Empiezan en \typ{0xFFFF}
o \typ{0xFF} para que la primera muestra siempre se imprima. No se
envían al teléfono.

\inosnippet{72}{84}

Estas líneas no son las funciones. Son el aviso al compilador de que
el cuerpo viene más abajo, para que \fnm{setup()} y \fnm{loop()}
puedan llamarlas. \fnm{serialInjectPoll} y
\fnm{handleSerialInjectLine} solo se declaran si
\nm{ENABLE\_SERIAL\_INJECT} vale 1. El cuerpo de cada una está en
la sección siguiente.

\inosnippet{87}{98}

\nm{ENABLE\_STEER\_SERVO} en 1 arma el servo. En 0 las dos
funciones del servo quedan vacías. \nm{PIN\_STEER\_SERVO} es el
GPIO 13. \nm{CH\_STEER\_SERVO} es el canal LEDC 2, el que usa el
núcleo 2.x. El \kw{\#ifndef} deja cambiar el valor desde el IDE
antes de este archivo, y si nadie lo cambió pone el de fábrica.

\inosnippet{100}{114}

\nm{kTelemetryChannels} tiene ocho filas y no admite una novena.
Cada llave es un \nm{TelChannel}: tipo, señal, pin, canal LEDC,
\nm{dutyMin} y \nm{dutyMax}. La fila 0 es CH1 (\str{v0}) y
la fila 7 es CH8 (\str{v7}).

De fábrica, las filas 0 a 6 son \nm{TEL\_VEHICLE}: temperatura,
acelerador, batería, dirección, luces, perilla izquierda y perilla
derecha. El pin y el canal van en $-1$ porque la señal no sale de un
GPIO. La fila 7 es \nm{TEL\_PWM} del servo: pin
\nm{PIN\_STEER\_SERVO}, canal \nm{CH\_STEER\_SERVO}, duty de
204 a 409 (el recorrido 1000--2000\,\textmu s).

\inosnippet{139}{140}

Estas dos constantes están justo antes de \fnm{steerDutyFromUs()}.
No están en la cabecera.
\nm{kSteerCenterUs} es 1500\,\textmu s, el centro.
\nm{kSteerTravelUs} es 500\,\textmu s a cada lado de la palanca.
\nm{rx} = 0 queda en 1500\,\textmu s.

\section{Qué hace cada función}
\label{sec:funcs}

El orden sigue el de cada archivo, después de las declaraciones. Debajo
de cada firma va la función entera, con el color del IDE. En el texto
se usa el mismo código: el naranja es una función, del sketch o de una
librería (\fnm{onControl()}, \fnm{print}, \fnm{ledcWrite}), y el morado
es un nombre de este proyecto (\nm{BT\_NAME}, \nm{g\_lx},
\nm{TEL\_TEMP}).

\nm{DBG\_PRINT}, \nm{DBG\_PRINTLN} y \nm{DBG\_PRINTF} no son
funciones. Con \nm{ENABLE\_DEBUG} en 1 escriben en el monitor serie.
En 0 el compilador las borra.

\subsection{En \texttt{TeleCon\_Classic\_Simple.ino}}

\begin{fnblock}
\fn{\typ{uint8\_t} \fnm{telemetryVehicle}(\nm{TelVehicle} \nm{id})}{116}{137}
Línea 116. Devuelve el byte 0--255 de una señal integrada, la que pide
una fila \nm{TEL\_VEHICLE}. \nm{TEL\_TEMP} es la constante
\nm{DEMO\_MOTOR\_TEMP}. \nm{TEL\_THROTTLE} es la magnitud de
\nm{g\_ly}. \nm{TEL\_BATT} escala \nm{DEMO\_BATT\_PCT} de
0--100 a 0--255. \nm{TEL\_STEER} es la magnitud de \nm{g\_rx}.
\nm{TEL\_LIGHTS} vale 255 si el bit 0 de \nm{g\_sw} está activo,
y 0 si no. \nm{TEL\_KNOB\_L} y \nm{TEL\_KNOB\_R} llevan la
perilla de 0--1023 a 0--255. \nm{TEL\_RY} es la magnitud de
\nm{g\_ry}. Otro \nm{id} devuelve 0. Acelerador, dirección y
\nm{TEL\_RY} no guardan el signo: los dos extremos de la palanca
dan el mismo byte.

\end{fnblock}

\begin{fnblock}
\fnstep{142}{145}{\typ{uint32\_t} \fnm{steerDutyFromUs}(\typ{int} \nm{us})}{%
Línea 142. Pasa un ancho de pulso, en microsegundos, a duty LEDC de
12 bits: pulso $\times 4095 / 20000$. Antes limita el pulso a
500--2500\,\textmu s. 1500\,\textmu s queda en el centro de un marco
de 20\,ms.}

\end{fnblock}

\begin{fnblock}
\fnstep{147}{151}{\typ{uint32\_t} \fnm{steerDutyFromRx}(\typ{int} \nm{rx})}{%
Línea 147. Limita \nm{rx} a $-100\ldots 100$, lo convierte en
microsegundos (0 es 1500\,\textmu s, a fondo $\pm 500$\,\textmu s) y
llama a \fnm{steerDutyFromUs}. El trim del teléfono ya viene dentro
de \nm{rx}.}

\end{fnblock}

\begin{fnblock}
\fntwo{\typ{void} \fnm{steerServoBegin}()}{153}{163}{174}{174}
Línea 154, o la línea 174 si \nm{ENABLE\_STEER\_SERVO} es 0.
Con el servo activo, deja un PWM de 50\,Hz y 12 bits en
\nm{PIN\_STEER\_SERVO} y escribe el centro. En el núcleo 3.x el
canal se une por número de pin (\fnm{ledcAttach},
\fnm{ledcWrite}). En el 2.x se usa \nm{CH\_STEER\_SERVO}
(\fnm{ledcSetup}, \fnm{ledcAttachPin}, \fnm{ledcWrite}).
Con el interruptor en 0 la función no hace nada.

\end{fnblock}

\begin{fnblock}
\fntwo{\typ{void} \fnm{steerServoWriteRx}(\typ{int} \nm{rx})}{165}{172}{175}{175}
Línea 165, o la línea 175 si el servo está apagado. Calcula el duty
de \nm{rx} y lo escribe. La versión vacía solo descarta
\nm{rx}, para que el compilador no avise de un parámetro sin uso.
\fnm{applyFailSafe()} la llama con 0, que es el centro.

\end{fnblock}

\begin{fnblock}
\fn{\typ{void} \fnm{onControl}(\typ{int} \nm{lx}, \typ{int} \nm{ly}, \typ{int} \nm{rx}, \typ{int} \nm{ry}, \typ{int} \nm{lk}, \typ{int} \nm{rk}, \typ{uint8\_t} \nm{sw})}{178}{187}
Línea 178. Corre una vez por cada línea \str{RC:CTRL}. Imprime el
cambio en el monitor, mueve el servo con \nm{rx}, y deja el bloque
\cmt{YOUR CODE HERE} para lo que deba seguir la palanca: PWM de
motores, un servo de paneo desde \nm{rk}, un pin de luces desde
\nm{sw} \& 0x01. Los casteos \texttt{(}\typ{void}\texttt{)} callan el aviso de
los valores que el arranque todavía no usa. Se borra el casteo del
valor que se empiece a usar.

\end{fnblock}

\begin{fnblock}
\fn{\typ{void} \fnm{onButton}(\typ{int} \nm{id})}{189}{198}
Línea 189. Corre una vez por cada \str{RC:BTN}. Imprime el id. El
comentario de fábrica deja el id 2 para un pulso corto de bocina y el
id 4 para centrar un servo de paneo. Esos dos cuerpos están vacíos. Un
botón que tenga que quedarse abajo todo el tiempo encaja mejor como
bit de \nm{sw} dentro de \fnm{onControl()}, porque
\str{RC:BTN} es la pulsación, no el nivel sostenido.

\end{fnblock}

\begin{fnblock}
\fn{\typ{void} \fnm{setup}()}{200}{242}
Línea 200. Arranca el serie a 115200, espera 200\,ms, centra el servo,
apaga Secure Simple Pairing, fija el PIN \str{1234} y abre
Bluetooth Classic con \nm{BT\_NAME}, como esclavo y sin BLE. Si
\nm{SerialBT}\texttt{.}\fnm{begin} falla, imprime el fallo y se queda en
\kw{while} (\typ{true}). El cartel de después es solo depuración:
nombre, saludo y líneas de ejemplo para el monitor.

\end{fnblock}

\begin{fnblock}
\fn{\typ{void} \fnm{loop}()}{244}{302}
Línea 244. Arduino llama a \fnm{loop()} una vez tras otra, desde que
termina \fnm{setup()}. Cada vuelta sigue estos pasos, y el siguiente
empieza cuando el anterior ya terminó.

\fnstep{245}{256}{1. Cliente Bluetooth}{%
  Lee \nm{SerialBT}\texttt{.}\fnm{hasClient()} y lo compara con
  \nm{btWasConnected}. Si cambió, guarda el estado nuevo.
  Si el teléfono acaba de entrar, llama a \fnm{resetHandshake()}:
  \nm{handshakeOk} queda en falso, y el monitor avisa que espera
  \str{RC:CONNECT,proto,simple}.
  Si acaba de salir, llama a \fnm{applyFailSafe()}
  (palancas en 0, perillas en 512, interruptores en 0, servo al
  centro) y después a \fnm{resetHandshake()}.
  Si el cliente sigue igual, este bloque no llama a nadie.}
\fnstep{258}{264}{2. Líneas de Bluetooth}{%
  Mientras \nm{SerialBT} tenga bytes, lee hasta
  \str{\textbackslash n}, quita espacios con \fnm{trim()} y, si queda
  texto, llama a \fnm{handleLine()}.
  Una línea \str{RC:CONNECT,proto,simple} válida en esta misma vuelta
  pone \nm{handshakeOk} en verdadero antes de los temporizadores de
  abajo.}
\fnstep{266}{268}{3. Monitor serie}{%
  Con \nm{ENABLE\_SERIAL\_INJECT} en 1, llama a
  \fnm{serialInjectPoll()}.
  Una línea del monitor que empiece por \str{RC:DATA} o por
  \str{RC:PLOT} sale al teléfono en ese momento.
  Con el interruptor en 0, este bloque no entra al binario.}
\fnstep{270}{270}{4. Compuerta del saludo}{%
  Si \nm{handshakeOk} es falso, \fnm{loop()} termina en ese
  \kw{return}.
  El tablero y el gráfico esperan la respuesta de
  \fnm{handleConnect()}: \str{RC:ACK,app,RC}.}
\fnstep{272}{272}{5. Reloj}{%
  Guarda \fnm{millis()} en \nm{now}.
  El tablero y el gráfico usan este mismo instante, cada uno con su
  marca. Cumplir un plazo no reinicia el otro.
  Los dos pueden cumplirse en la misma vuelta: primero el tablero,
  después el gráfico.}
\fnstep{274}{286}{6. Tablero, cada \nm{TELEMETRY\_DATA\_MS} (500\,ms)}{%
  Cuando \nm{now} $-$ \nm{lastDataMs} llega al plazo, guarda
  \nm{now} en \nm{lastDataMs} y arma una línea \str{RC:DATA}.
  \str{left} es \fnm{abs}(\nm{g\_ly}) * 25 / 10: acelerador 100
  sale como 250, y el teléfono lo lee como velocidad $\times 10$.
  \fnm{abs()} manda la magnitud, así que acelerador a fondo hacia
  atrás también sale como 250.
  \str{right} va fijo en 0.
  \str{lo,1,ro,0,lg,1,rg,0} son las lámparas del panel, también fijas.
  \str{analog} copia \nm{DEMO\_MOTOR\_TEMP} (90) y \str{batt} copia
  \nm{DEMO\_BATT\_PCT} (76).
  \str{led} es \nm{g\_sw} en dos dígitos hexadecimales.
  \fnm{sendLine()} manda esa línea.}
\fnstep{288}{301}{7. Gráfico, cada \nm{PLOT\_MS} (50\,ms)}{%
  Este bloque entra al binario cuando \nm{ENABLE\_PLOT} vale 1.
  Cuando \nm{now} $-$ \nm{lastPlotMs} llega al plazo, guarda
  \nm{now} en \nm{lastPlotMs}, llama a \fnm{telemetryReadAll()} con
  \nm{kTelemetryChannels} y \fnm{telemetryVehicle}, y manda
  \str{RC:PLOT} de \str{v0} a \str{v7}, en el orden de las ocho filas.
  Con \nm{ENABLE\_PLOT} en 0, este bloque no entra al binario.}

\end{fnblock}

\begin{fnblock}
\fn{\typ{void} \fnm{handleSerialInjectLine}(\typ{const} \typ{String}\& \nm{line})}{304}{324}
Línea 305. El sketch de fábrica la compila, porque
\nm{ENABLE\_SERIAL\_INJECT} vale 1. Si no hay cliente Bluetooth, o
el saludo todavía no está listo, no envía la línea. Si empieza por
\str{RC:DATA} o por \str{RC:PLOT}, la manda una sola vez con
\fnm{sendLine}. Otra línea solo se anota en el monitor. Los
temporizadores de \fnm{loop()} siguen, así el paquete automático
siguiente reemplaza esa línea en el teléfono.

\end{fnblock}

\begin{fnblock}
\fn{\typ{void} \fnm{serialInjectPoll}()}{326}{347}
Línea 326. Lee el monitor carácter a carácter hasta \str{\textbackslash n}.
Ignora \str{\textbackslash r}, para que sirva un fin de línea
``Both NL \& CR''. Al cerrar la línea la recorta y, si queda texto,
llama a \fnm{handleSerialInjectLine}. Si el buffer pasa de
\nm{SERIAL\_MAX\_LINE} (192), lo tira y avisa.

\end{fnblock}

\begin{fnblock}
\fnstep{349}{351}{\typ{uint16\_t} \fnm{stickTo12}(\typ{int} \nm{s})}{%
Línea 349. Lleva una palanca de $-100\ldots 100$ a $0\ldots 4095$, solo
para el texto del monitor. 0 en la palanca queda cerca de 2048. Ese
número no es el duty del servo ni el byte de \str{RC:PLOT}.}

\end{fnblock}

\begin{fnblock}
\fn{\typ{void} \fnm{debugPrintControl}(\typ{int} \nm{lx}, \typ{int} \nm{ly}, \typ{int} \nm{rx}, \typ{int} \nm{ry}, \typ{int} \nm{lk}, \typ{int} \nm{rk}, \typ{uint8\_t} \nm{sw})}{353}{398}
Línea 353. Pasa cada palanca por \fnm{stickTo12} y cada perilla por
un límite 0--1023. Imprime un eje solo cuando ese valor cambió respecto
de \nm{lastPrint*}, así una palanca quieta no llena el monitor.
Si cambia uno de los seis primeros bits de \nm{sw}, imprime
\str{SW1} a \str{SW6}. Con \nm{ENABLE\_DEBUG} en 0 la función
sigue corriendo y las macros no escriben nada.

\end{fnblock}

\begin{fnblock}
\fnstep{400}{402}{\typ{void} \fnm{debugPrintButton}(\typ{int} \nm{id})}{%
Línea 400. Imprime \str{BTNn : Pressed} con el id que llegó en
\str{RC:BTN}.}

\end{fnblock}

\begin{fnblock}
\fnstep{404}{407}{\typ{void} \fnm{sendLine}(\typ{const} \typ{String}\& \nm{line})}{%
Línea 404. Único envío hacia el teléfono. Escribe el texto y después
un salto de línea, por \nm{SerialBT}. El saludo, el NAK, el tablero
y el gráfico pasan todos por aquí.}

\end{fnblock}

\begin{fnblock}
\fnstep{409}{416}{\typ{void} \fnm{applyFailSafe}()}{%
Línea 409. Deja las cuatro palancas en 0, las dos perillas en 512 y
\nm{g\_sw} en 0, y escribe el servo con \nm{rx} 0 (centro).
Debajo del segundo \cmt{YOUR CODE HERE} se paran los motores que
\fnm{onControl()} haya puesto en marcha. Sin eso, el último
acelerador sigue aplicado cuando el teléfono ya no está.}

\end{fnblock}

\begin{fnblock}
\fnstep{418}{423}{\typ{void} \fnm{resetHandshake}()}{%
Línea 418. Pone \nm{handshakeOk} en falso y olvida los últimos
valores del monitor, para que la sesión nueva vuelva a imprimirlos.
No mueve el servo ni las salidas: quien vio la caída ya llamó a
\fnm{applyFailSafe()}.}

\end{fnblock}

\begin{fnblock}
\fn{\typ{String} \fnm{getValue}(\typ{const} \typ{String}\& \nm{line}, \typ{const} \typ{char}* \nm{key})}{425}{447}
Línea 425. Saca el valor de una clave en una línea
\str{APP:TIPO,clave,valor,clave,valor}. Salta la cabeza hasta la
primera coma y recorre los pares. El valor llega hasta la coma
siguiente, o hasta el fin de la línea en el último par. Las claves
distinguen mayúsculas (\nm{lx}, no \str{LX}). Si la clave no
está, devuelve una cadena vacía, y \fnm{toInt()} de esa cadena es 0.

\end{fnblock}

\begin{fnblock}
\fn{\typ{void} \fnm{handleConnect}(\typ{const} \typ{String}\& \nm{app}, \typ{const} \typ{String}\& \nm{line})}{449}{475}
Línea 449. Exige dos cosas, sin distinguir mayúsculas. El nombre antes
de los dos puntos tiene que ser \nm{APP\_PREFIX} (\str{RC}). El
valor \str{proto} tiene que ser \nm{PROTO\_WIRE}
(\str{simple}). Si el nombre no coincide, responde
\str{RC:NAK} con \str{app\_mismatch} y deja
\nm{handshakeOk} en falso. Si falta \str{proto}, o trae otro
valor, responde \str{proto\_mismatch}. Si ambos coinciden, envía
exactamente \str{RC:ACK,app,RC} y pone \nm{handshakeOk} en
verdadero.

\end{fnblock}

\begin{fnblock}
\fn{\typ{void} \fnm{handleLine}(\typ{const} \typ{String}\& \nm{line})}{477}{508}
Línea 477. Si la línea no tiene dos puntos, vuelve. Separa la app y el
tipo (el texto hasta la primera coma). \str{CONNECT} llama a
\fnm{handleConnect} aunque el saludo no esté listo. Cualquier otro
tipo se ignora mientras \nm{handshakeOk} es falso, y también si la
app no es \str{RC}. \str{CTRL} guarda \nm{lx}, \nm{ly},
\nm{rx}, \nm{ry}, \nm{lk}, \nm{rk} y \nm{sw} (el
campo \nm{sw} se lee en hexadecimal) y llama a \fnm{onControl}.
\str{BTN} llama a \fnm{onButton} con el \nm{id}. Una línea
\str{DATA}, \str{PLOT} o \str{SET} que llegue del teléfono no
hace nada: este sketch las envía, no las obedece.

\end{fnblock}

\subsection{En \texttt{TelemetryChannels.h}}

El archivo declara los tipos que usan las tres funciones.
\nm{TelKind} es el origen de la fila: vehículo, ADC, GPIO o PWM.
\nm{TelVehicle} es la señal integrada (\nm{TEL\_TEMP},
\nm{TEL\_THROTTLE}, \nm{TEL\_BATT}, \nm{TEL\_STEER},
\nm{TEL\_LIGHTS}, \nm{TEL\_KNOB\_L}, \nm{TEL\_KNOB\_R},
\nm{TEL\_RY}). \nm{TelChannel} guarda el tipo, la señal, el
pin, el canal LEDC y el rango de duty. \nm{TelVehicleFn} es el
puntero a \fnm{telemetryVehicle()}.

\begin{fnblock}
\fnh{\typ{uint8\_t} \fnm{telemetryReadPwm}(\typ{const} \nm{TelChannel}\& \nm{ch})}{46}{56}
Línea 46. Si \nm{dutyMax} no es mayor que \nm{dutyMin},
devuelve 0. Lee el duty que ya está corriendo: en el núcleo 3.x con
\fnm{ledcRead} del número de pin, en el 2.x con \fnm{ledcRead}
del número de canal. Recorta ese duty a \nm{dutyMin}\ldots
\nm{dutyMax} y lo escala a 0--255. La fila de fábrica del servo
usa 204\ldots409, que es el recorrido 1000--2000\,\textmu s; el centro
sale cerca de 128. El pin tiene que estar unido en \fnm{setup}
antes de que esta función lo lea.

\end{fnblock}

\begin{fnblock}
\fnh{\typ{uint8\_t} \fnm{telemetryReadOne}(\typ{const} \nm{TelChannel}\& \nm{ch}, \nm{TelVehicleFn} \nm{vehicle})}{58}{71}
Línea 58. Un byte, según \nm{ch}\texttt{.}\nm{kind}. \nm{TEL\_ADC} lee el
pin y escala 0--4095 a 0--255; un pin negativo devuelve 0.
\nm{TEL\_GPIO} devuelve 255 o 0 según \fnm{digitalRead}; un pin
de salida informa el nivel que se escribió. \nm{TEL\_PWM} llama a
\fnm{telemetryReadPwm}. Cualquier otro tipo llama a
\fnm{vehicle}(\nm{ch}\texttt{.}\nm{vehicle}). Si ese puntero es nulo, devuelve 0.

\end{fnblock}

\begin{fnblock}
\fnhstep{73}{75}{\typ{void} \fnm{telemetryReadAll}()}{%
Línea 73. Recorre las ocho filas de \nm{kTelemetryChannels[]} y
escribe \nm{out[0]} hasta \nm{out[7]}. \fnm{loop()} manda
esos bytes como \str{v0} hasta \str{v7}. La fila 0 es CH1 del
teléfono y la fila 7 es CH8. No hay noveno puesto en esta función.}


\end{fnblock}

\appendix
\renewcommand{\thesection}{Anexo~\Alph{section}}
\section{\texttt{TeleCon\_Classic\_Simple.ino}}
\label{sec:ino-full}

Este es el sketch entero, el archivo que abre el IDE. Los números de
línea son los del archivo.

\lstinputlisting[basicstyle=\ttfamily\footnotesize,%
  numberstyle=\scriptsize\color{ArduinoLineNo}]{\sketchfile}

\section{\texttt{TelemetryChannels.h}}
\label{sec:h-full}

La cabecera va aparte. El \texttt{.ino} la incluye y no la copia.

\lstinputlisting[basicstyle=\ttfamily\footnotesize,%
  numberstyle=\scriptsize\color{ArduinoLineNo}]{\channelfile}

\end{document}
